\subsection{Depth and Koszul complex}

Let A be a ring, M an A-module, \(\mathbf{x} = (x_i)_{i \in I}\) a family of elements of A. Denote by \(u : \mathrm{A}^{(1)} \to \mathrm{A}\) the linear form such that \(u(c_i) = x_i\) for every \(i \in I\) , and by \(\mathbf{K}^{\bullet}(\mathbf{x}, \mathrm{M})\) the complex \(\mathbf{K}_{\mathrm{A}}^{\bullet}(u, \mathrm{M})\) associated with u (A, X, p.~147). One has \(\mathbf{K}^p(\mathbf{x}, \mathrm{M}) = 0\) for p \textless{} 0 ; for \(p \geqslant 0\) the A-module \(\mathbf{K}^p(\mathbf{x}, \mathrm{M}) = \operatorname{Hom}_{\Lambda} (\boldsymbol{\Lambda}^p(\mathrm{A}^{(1)}), \mathrm{M})\) is canonically identified with the A-module \(C_I^p(\mathrm{M})\) formed by the alternating maps from \(I^p\) to M (A, X, p.~153), the differential \(\partial^p : \mathbf{K}^p(\mathbf{x}, \mathrm{M}) \longrightarrow \mathbf{K}^{p+1}(\mathbf{x}, \mathrm{M})\) being given by the formula

\[
\left(\partial^ {p} \boldsymbol {m}\right) \left(\alpha_ {1}, \dots , \alpha_ {p + 1}\right) = \sum_ {j = 1} ^ {p + 1} (- 1) ^ {j + 1} x _ {\alpha_ {j}} \boldsymbol {m} \left(\alpha_ {1}, \dots , \alpha_ {j - 1}, \alpha_ {j + 1}, \dots , \alpha_ {p + 1}\right)
\]

for \(\boldsymbol{m} \in \mathbf{K}^{p}(\mathbf{x}, \mathrm{M})\) and \((\alpha_{1}, \ldots, \alpha_{p+1}) \in \mathrm{I}^{p+1}\) (A, X, p. 154, formula (12)). It follows in particular that the complex \(\mathbf{K}^{\bullet}(\mathbf{x}, \mathrm{M})\) depends only on the structure of M as a Z-module and on the endomorphisms \((x_{i})_{\mathrm{M}}\) .

One denotes by \(\mathrm{H}^{\bullet}(\mathbf{x},\mathrm{M})\) the homology of the complex \(\mathbf{K}^{\bullet}(\mathbf{x},\mathrm{M})\) . The A-module \(\mathrm{H}^{0}(\mathbf{x},\mathrm{M})\) is identified with \(\mathrm{Hom}_{\Lambda}(\mathrm{A}/\mathrm{J},\mathrm{M})\) , where J is the ideal of A generated by the \(x_{i}\) (A, X, p.~147, lemma 1).

Let \((\mathrm{M}_{\alpha})_{\alpha\in\mathrm{K}}\) be a family of A-modules, and M its product; the complex \(\mathbf{K}^{\bullet}(\mathbf{x},\mathrm{M})\) is canonically isomorphic to the product complex of the \(\mathbf{K}^{\bullet}(\mathbf{x},\mathrm{M}_{\alpha})\) , so that for each integer s the A-module \(\mathrm{H}^{s}(\mathbf{x},\mathrm{M})\) is identified with the product of the \(\mathrm{H}^{s}(\mathbf{x},\mathrm{M}_{\alpha})\) (A, X, p. 28, prop. 1).

THEOREM 1.-- Let A be a ring, J an ideal of A, \(\mathbf{x} = (x_i)_{i \in I}\) a generating family of J, M an A-module. The depth of M relative to J is the lower bound (in \(\mathbf{N} \cup \{+\infty\}\) ) of the integers \(n\) such that \(\Pi^n(\mathbf{x}, \mathbf{M}) \neq 0\) .

Set \(p = \text{prof}_{\Lambda}(\text{J}; \text{M})\) . Consider the complex \(\mathbf{K}^{\bullet}(\mathbf{x}, \text{M})\) . Its homology is annihilated by J ( \(\Lambda\) , X, p.~148, cor. 2), and the depth relative to J of each of the modules \(\mathbf{K}^{i}(\mathbf{x}, \text{M})\) is equal to \(p\) or to \(+\infty\) (no. 1, remark 4). It then follows from cor. 1 of no. 2 that one has \(\text{H}^{i}(\text{x}, \text{M}) = 0\) for \(i < p\) . It remains to prove that \(\text{H}^{p}(\text{x}, \text{M})\) is nonzero when \(p < +\infty\) .

The case \(p = 0\) being evident, suppose \(0 < p < +\infty\) and \(\mathrm{H}^p(\mathbf{x},\mathrm{M}) = 0\) . Let L be a free resolution of the A-module \(\mathrm{A / J}\) ; denote by C the complex \(\mathrm{Homgr}_{\mathrm{A}}(\mathrm{L},\mathrm{M})\) . The \(\Lambda\) -module \(\mathrm{H}^i(\mathrm{C})\) is isomorphic to \(\mathrm{Ext}_{\mathrm{A}}^i(\mathrm{A / J},\mathrm{M})\) (A, X, p.~100, th. 1); it is therefore zero for i \textless{} p.~One then has for i \textless{} p canonical exact sequences

\[
0 \rightarrow \mathrm{B} ^ {i} (\mathrm{C}) \rightarrow \mathrm{C} ^ {i} \rightarrow \mathrm{B} ^ {i + 1} (\mathrm{C}) \rightarrow 0.
\]

The A-module \(C^{i}\) is a product of A-modules isomorphic to M ; hence one has \(\mathrm{H}^{s}(\mathbf{x},\mathrm{C}^{i})=0\) for \(s\leqslant p\) . From the preceding exact sequences and A, X, p.~150 one deduces that the connecting homomorphism \(\partial^{s}:\mathrm{H}^{s}(\mathbf{x},\mathrm{B}^{i+1}(\mathrm{C}))\longrightarrow\mathrm{H}^{s+1}(\mathbf{x},\mathrm{B}^{i}(\mathrm{C}))\) is injective for \(s\leqslant p\) and i\textless p ; since \(\mathrm{B}^{0}(\mathrm{C})=0\) , it follows that \(\mathrm{H}^{p-i}(\mathbf{x},\mathrm{B}^{i+1}(\mathrm{C}))\) is zero for i\textless p.~In particular, we have \(\mathrm{II}^{1}(\mathbf{x},\mathrm{B}^{p}(\mathrm{C}))=0\) , so that the exact sequence

\[
0 \to \mathrm{B} ^ {p} (\mathrm{C}) \to \mathrm{Z} ^ {p} (\mathrm{C}) \to \mathrm{H} ^ {p} (\mathrm{C}) \to 0
\]

provides a surjection \(\mathrm{H}^{0}(\mathbf{x},\mathrm{Z}^{p}(\mathrm{C}))\longrightarrow\mathrm{H}^{0}(\mathbf{x},\mathrm{H}^{p}(\mathrm{C}))\) . Since \(\mathrm{H}^{p}(\mathrm{C})\) is isomorphic to \(\operatorname{Ext}_{\mathrm{A}}^{p}(\mathrm{A}/\mathrm{J},\mathrm{M})\) which is nonzero and annihilated by J, we have \(\mathrm{H}^{0}(\mathbf{x},\mathrm{H}^{p}(\mathrm{C}))\neq0\) , whence \(\mathrm{H}^{0}(\mathbf{x},\mathrm{Z}^{p}(\mathrm{C}))\neq0\) and consequently \(\mathrm{H}^{0}(\mathbf{x},\mathrm{C}^{p})\neq0\) . But this implies \(\mathrm{H}^{0}(\mathbf{x},\mathrm{M})\neq0\) , contrary to the hypothesis. We therefore have \(\mathrm{H}^{p}(\mathbf{x},\mathrm{M})\neq0\) , which completes the proof.

COROLLARY 1.--- Suppose the ideal J is finitely generated and JM ≠ M. Then prof \(_{A}\) (J; M) is finite and ≤ Card(I); for it to be equal to Card(I), it is necessary and sufficient that the family x be completely secant for M (A, X, p. 157, def. 2).

Suppose first that 1 is finite, and denote \(r\) its cardinality. The A-module \(\mathrm{H}^r (\mathbf{x},\mathrm{M})\) is canonically isomorphic to \(\mathrm{H}_0(\mathbf{x},\mathrm{M})\) , itself isomorphic to M/JM (A, X, p.~155); the inequality \(\mathrm{prof}_{\Lambda}(\mathrm{J};\mathrm{M})\leqslant r\) It therefore follows from Th. 1. For equality to hold, it is necessary and sufficient that the A-module \(\mathrm{H}^i (\mathbf{x},\mathrm{M})\) be zero for \(i < r\) , which means that the family \(\mathbf{x}\) is completely secant for M (A, X, p.~157).

Based on what precedes, \(\mathrm{prof}_{\mathrm{A}}(\mathrm{J};\mathrm{M})\) is finite; it remains to prove that if the family \(\mathbf{x}\) is completely secant for M, the set I is finite. But the condition \(\mathrm{H}_{1}(\mathbf{x},\mathrm{M}) = 0\) (A, X, p. 157, def. 2) implies that we have an exact sequence

\[
\boldsymbol {\Lambda} ^ {2} (\mathrm{A} ^ {(\mathrm{I})}) \otimes_ {\mathrm{A}} \mathrm{M} \xrightarrow {\partial_ {2}} \mathrm{M} ^ {(\mathrm{I})} \xrightarrow {\partial_ {1}} \mathrm{JM} \to 0,
\]

where the image of \(\partial_{2}\) is contained in \(\mathrm{JM}^{(1)}\) By taking the tensor product with A/J, we deduce an A-linear isomorphism from \((\mathrm{M} / \mathrm{JM})^{(1)}\) on \(\mathrm{JM} / \mathrm{J}^2\mathrm{M}\) . Now this latter module is of finite type, since J and M are; as M/JM is not zero, it follows that the set I is finite.

COROLLARY 2. Let A be a local ring, J a finitely generated ideal of A distinct from A, M a nonzero finitely generated A-module. Set \(r = [\mathrm{J} / \mathfrak{m}_{\mathrm{A}}\mathrm{J} : \kappa_{\mathrm{A}}]\) . One has \(\operatorname{prof}_{\Lambda}(\mathrm{J};\mathrm{M}) \leqslant r\) ; there is equality if and only if J is generated by a completely secant family for M. In this case, for a generating family of J to be completely secant, it is necessary and sufficient that it have r elements.

By Nakayama's lemma, we have \(\mathrm{JM} \neq \mathrm{M}\) , and \(r\) is the minimal number of generators of \(\mathbf{J}\) ; cor. 2 therefore follows from cor. 1.

PROPOSITION 4.--- Let \(\rho: A \to B\) a homomorphism of rings, \(J\) an ideal of \(A\) and \(N\) a B-module. We have the equality \(\operatorname{prof}_{\mathrm{A}}(\mathrm{J}; \mathrm{N}) = \operatorname{prof}_{\mathrm{B}}(\mathrm{JB}; \mathrm{N})\) .

Let \(\mathbf{x} = (x_{i})_{i \in \mathbb{I}}\) a generating family of J; the family \(\rho(\mathbf{x}) = (\rho(x_{i}))_{i \in \mathbb{I}}\) generates JB. By construction the complex \(\mathbf{K}^{\bullet}(\boldsymbol{\rho}(\mathbf{x}), \mathbf{N})\) is equal to \(\mathbf{K}^{\bullet}(\mathbf{x}, \mathbf{N})\) . The proposition therefore follows from th. 1.

COROLLARY. Let A be a local ring, a an ideal of A distinct from A and M an A-module annihilated by a. We have \(\operatorname{prof}_{\mathrm{A}}(\mathrm{M}) = \operatorname{prof}_{\mathrm{A}/\mathfrak{a}}(\mathrm{M})\) .

Let \(\rho: A \to B\) a homomorphism of rings, \(x = (x_i)_{i \in I}\) a finite family of elements of \(\Lambda\) and M an A-module. For every integer \(p\) let us denote \(u^p: B \otimes_\Lambda C_I^p(M) \longrightarrow C_I^p(B \otimes_\Lambda M)\) the B-linear homomorphism which associates to \(b \otimes m\) the alternating map \((\alpha_1, \ldots, \alpha_p) \mapsto b \otimes m(\alpha_1, \ldots, \alpha_p)\) . The family \((u^p)\) defines an isomorphism of complexes

\[
u: \mathrm{B} \otimes_ {\mathrm{A}} \mathbf {K} ^ {\bullet} (\mathbf {x}, \mathrm{M}) \longrightarrow \mathbf {K} ^ {\bullet} (\mathbf {x}, \mathrm{B} \otimes_ {\mathrm{A}} \mathrm{M}).
\]

Consider the canonical homomorphism

\[
\gamma^ {p} (\mathrm{B}, \mathbf {K} ^ {\bullet} (\mathbf {x}, \mathrm{M})): \mathrm{B} \otimes_ {\mathrm{A}} \mathrm{H} ^ {p} (\mathbf {x}, \mathrm{M}) \longrightarrow \mathrm{H} ^ {p} (\mathrm{B} \otimes_ {\mathrm{A}} \mathbf {K} ^ {\bullet} (\mathbf {x}, \mathrm{M}))
\]

(A, X, p. 62); by composition with \(\mathrm{H}^{p}(u)\) , we deduce from it a homomorphism

\[
v ^ {p}: \mathrm{B} \otimes_ {\Lambda} \mathrm{H} ^ {p} (\mathbf {x}, \mathrm{M}) \longrightarrow \mathrm{H} ^ {p} (\mathbf {x}, \mathrm{B} \otimes_ {\Lambda} \mathrm{M}).
\]

Lemma 1.-- If the A-module B is flat, the homomorphism \(v^p\) is bijective for every integer \(p\) .

This follows from A, X, p. 66, cor. 2.

PROPOSITION 5.--- Let A be a ring, J a finitely generated ideal of A and M an A-module. Let Ω denote the set of maximal ideals of A belonging to Supp(M) and containing J. We have

\[
\operatorname{prof} _ {A} (J; M) = \inf _ {\mathfrak {p} \in \operatorname{Spec} (A)} \operatorname{prof} _ {A _ {\mathfrak {p}}} \left(J _ {\mathfrak {p}}; M _ {\mathfrak {p}}\right) = \inf _ {m \in \Omega} \operatorname{prof} _ {\Lambda_ {m}} \left(J _ {m}; M _ {m}\right).
\]

Let \(\mathbf{x} = (x_{i})_{i \in 1}\) a finite generating family of J. Let p be a prime ideal of A; the ideal \(J_{p}\) is generated by the image \(x_{p}\) of the family \((x_{i})\) in \(A_{p}\) . For every \(p \geqslant 0\) , the \(A_{p}\) -module \(\left(\mathrm{H}^{p}(\mathbf{x},\mathrm{M})\right)_{\mathfrak{p}}\) is isomorphic to \(\mathrm{H}^{p}(\mathbf{x}_{\mathfrak{p}},\mathrm{M}_{\mathfrak{p}})\) (lemma 1); by th. 1, we therefore have \(\operatorname{prof}_{\Lambda}(\mathrm{J};\mathrm{M}) \leqslant \inf_{\mathfrak{p} \in \operatorname{Spec}(\mathrm{A})} \operatorname{prof}_{\mathrm{A}_{\mathfrak{p}}}(\mathrm{J}_{\mathfrak{p}}; \mathrm{M}_{\mathfrak{p}}) \leqslant \inf_{\mathfrak{m} \in \Omega} \operatorname{prof}_{\mathrm{A}_{\mathfrak{m}}}(\mathrm{J}_{\mathfrak{m}}; \mathrm{M}_{\mathfrak{m}})\) .

Let \(p\) be an integer strictly less than \(\operatorname{prof}_{\mathrm{A}_{\mathfrak{m}}}\left(\mathrm{JA}_{\mathfrak{m}};\mathrm{M}_{\mathfrak{m}}\right)\) for every \(\mathfrak{m}\in\Omega\) . One then has \(\mathrm{H}^{p}\left(\mathbf{x}_{\mathfrak{m}},\mathrm{M}_{\mathfrak{m}}\right)=0\) for every maximal ideal \(\mathfrak{m}\) of A: this follows from th. 1 if \(\mathfrak{m}\in\Omega\) , from the fact that \(\mathrm{M}_{\mathfrak{m}}=0\) if \(\mathfrak{m}\notin\operatorname{Supp}(\mathrm{M})\) , and from the fact that the ideal \(\mathrm{JA}_{\mathfrak{m}}\) , which annihilates \(\mathrm{H}^{p}\left(\mathbf{x}_{\mathfrak{m}},\mathrm{M}_{\mathfrak{m}}\right)\) (A, X, p. 148, cor. 2), is equal to \(\mathrm{A}_{\mathfrak{m}}\) if \(\mathfrak{m}\notin\mathrm{V}(\mathrm{J})\) . One therefore has \(\left(\mathrm{H}^{p}(\mathbf{x},\mathrm{M})\right)_{\mathfrak{m}}=0\) for every maximal ideal \(\mathfrak{m}\) of A, which implies \(\mathrm{H}^{p}(\mathbf{x},\mathrm{M})=0\) (II, § 3, n° 3, cor. 2 of th. 1). The proposition then follows from th. 1.

PROPOSITION 6.--- Let A be a ring, J a finitely generated ideal of A and M an A-module. Let B be a ring and \(\rho: A \to B\) a ring homomorphism making B into a flat A-module.

\begin{enumerate}
\def\labelenumi{\alph{enumi})}
\item
  One \(a \operatorname{prof}_{\mathrm{A}}(\mathrm{J}; \mathrm{M}) \leqslant \operatorname{prof}_{\mathrm{B}}(\mathrm{JB}; \mathrm{B} \otimes_{\mathrm{A}} \mathrm{M})\) .
\item
  Suppose moreover that every maximal ideal of Supp(M) containing J belongs to the image of the canonical map \(\operatorname{Spec}(B) \to \operatorname{Spec}(A)\) . One then has \(\operatorname{prof}_{\mathrm{A}}(\mathrm{J}; \mathrm{M}) = \operatorname{prof}_{\mathrm{B}}(\mathrm{JB}; \mathrm{B} \otimes_{\mathrm{A}} \mathrm{M})\) . This is the case in particular if the A-module B is faithfully flat.
\end{enumerate}

Assertion a) follows from th. 1 and lemma 1.

Let \(p\) be an integer strictly less than \(\text{prof}_B(\text{JB};\text{B} \otimes_A \text{M})\) , and let \(\mathfrak{m}\) be a maximal ideal of A belonging to \(\text{Supp(M)} \cap \text{V(J)}\) . Let \(\mathbf{x}\) be a finite generating family of the ideal J. Under the hypothesis of b), there exists a prime ideal \(\mathfrak{n}\) of B lying over \(\mathfrak{m}\) , and one has a canonical isomorphism

\[
\mathrm{B} _ {\mathfrak {n}} \otimes_ {\mathrm{A} _ {\mathfrak {m}}} \left(\mathrm{A} _ {\mathfrak {m}} \otimes_ {\Lambda} \mathrm{H} ^ {p} (\mathbf {x}, \mathrm{M})\right) \longrightarrow \mathrm{B} _ {\mathfrak {n}} \otimes_ {\mathrm{B}} \left(\mathrm{B} \otimes_ {\mathrm{A}} \mathrm{H} ^ {p} (\mathbf {x}, \mathrm{M})\right).
\]

Now \(\mathrm{B}\otimes_{\Lambda}\mathrm{H}^{p}(\mathbf{x},\mathrm{M})\) is isomorphic to \(\mathrm{H}^{p}(\mathsf{p}(\mathbf{x}),\mathrm{B}\otimes_{\Lambda}\mathrm{M})\) (lemma 1), hence is zero; moreover \(B_{n}\) is faithfully flat over \(A_{m}\) (I, § 3, n° 5, prop. 9 and II, § 3, n° 4, prop. 14 and remark). One therefore has \(A_{m}\otimes_{\Lambda}\mathrm{H}^{p}(\mathbf{x},\mathrm{M})=0\) and consequently p \textless{} prof \(_{A_{m}}(J_{m};M_{m})\) (lemma 1 and th. 1). The first assertion of b) then follows from prop. 5; the second follows from I, § 3, n° 5, prop. 9).

COROLLARY.--- Let A be a Noetherian ring, J an ideal of A, M a finitely generated A-module, \(\widehat{\Lambda}\) and \(\widehat{M}\) the separated completions of A and M for the J-adic topology. Then one has \(\operatorname{prof}_{\mathbf{A}}(\mathbf{J};\mathbf{M}) = \operatorname{prof}_{\widehat{\mathbf{A}}}(\widehat{\mathbf{J}}\widehat{\mathbf{A}};\widehat{\mathbf{M}})\) .

Indeed, the A-module \(\widehat{\mathsf{A}}\) is flat and the \(\widehat{\mathsf{A}}\) -module \(\widehat{\mathsf{M}}\) is isomorphic to \(\widehat{\mathsf{A}}\otimes_{\mathsf{A}}\mathsf{M}\) (III, § 3, n° 4, th. 3); moreover, every maximal ideal of A containing J belongs to the image of the map \(\operatorname{Spec}(\widehat{\mathsf{A}})\to\operatorname{Spec}(\mathsf{A})\) (loc. cit., prop. 8).

